% Appendix E -- the engineering lessons behind the 37.0 -> 12.3 ms reduction,
% moved out of section 5 (2026-09-26) for the page limit. The body keeps the
% number the model uses and what it is made of; this is how it got there, which
% is useful to someone rebuilding the path and not to someone reading the
% argument.

\section{Lessons from the transfer path}
\label{app:impl}

Beyond the transfer itself, using the host path costs a fixed \textbf{12.3\,ms}
per step, reduced from 37.0. That residual is the $F$ of
\S\ref{sec:model:omits}, on which the paper's result turns, so how it was found
and what it is made of both matter.

\paragraph{Layout decides whether the DMA engine can do the gather.} The mirrors
are block-major with $K$ and $V$ in separate tensors, so a run of consecutive
host blocks is one contiguous byte range on both sides and a layer's transfer is
a few \texttt{cudaMemcpyAsync} calls from pinned memory. An earlier layout
interleaved $K$ and $V$, making a run two chunks strided by a whole plane: no
pinned DMA was possible, and the resulting host-side gather cost 22\,ms per
layer, 98\% of the streaming path. The general point is that this cost is
invisible to the metric one would naturally watch, because it is a \emph{host}
memcpy that a faster interconnect does not touch --- on a coherent link the same
code would spend hundreds of milliseconds per step gathering in order to hide a
transfer of a few, and the natural conclusion would have been that the
interconnect does not help.

\paragraph{Hoist schedule-invariant work out of the per-layer path.} Four sites
executed work inside the per-layer loop that is fixed once the step's schedule is
known: a device reduction and a \texttt{.item()} inside a guard whose purpose was
to \emph{skip} work, scatter indices rebuilt per layer, a sequence maximum and an
\texttt{arange}, and a clone and index-fill of the masked sequence lengths.
Together they were worth 21\,ms per step, moving the end-to-end throughput
deficit from $-11.3\%$ to $-2.3\%$.

\paragraph{Time asynchronous GPU work with operator diffs, not phase timers.}
Aggregate per-phase timers misattributed this cost twice, plausibly each time: a
timer wrapped around asynchronous work charges waiting-for-the-GPU to whichever
phase holds it, and a profiler that synchronises inside its own timing wrapper
makes double-buffering impossible by construction and then reports that
double-buffering does not help. What located all four sites was an
operator-level trace differencing a step with one host-tier request against a
step with none. The same discipline applies to absolutes: the prefill handover of \S\ref{sec:impl:stash}
first measured 4.06\,GB/s on a contended node against 38\,GB/s exclusive on an
idle one, a $9.4\times$ error from contention, which is why
\S\ref{sec:eval:setup} requires exclusive access for every absolute.

\paragraph{What the residual is made of.} An operator trace shows
\texttt{cudaStreamSynchronize} gone and no dominant remaining term: roughly 640
extra kernel launches per step, about ten per layer across 64 layers, costing
$0.19$\,ms \emph{per layer}, so $64 \times 0.19 = 12.2$\,ms --- the per-layer structure itself. Removing it needs batching
across layers rather than issuing per layer, which we have not built. It is
launch-bound rather than bandwidth-bound, so a faster interconnect does not
shrink it; but it is \emph{not} machine-independent, since launch cost depends on
the host CPU and the count scales with layer depth. Whether it is smaller on a
Grace host is the first thing \S\ref{sec:eval} must measure.
